% ===== §17 The Implementation-Level Interface =====
\section{Observable Signatures: An Implementation-Level Interface}
\label{p22:sec:neuro}

\noindent
The account so far has been ontological, ethical, and cultivational, and it would be incomplete in a particular way if it stopped there, because it has repeatedly made claims of a kind that ought to leave traces. If the perception of the field is the relaxation of an over-weighted prior, if reception is the integration of what was met rather than its defence, if sharing is a coupling event between two subjects, and if reproduction is a threshold lowered across cycles, then each of these, were it actual, should show up in something measurable, and a framework that named such mechanisms while refusing every commitment about their signatures would be protecting itself from test in a way the series has elsewhere refused to. This section supplies the interface: for each of the four movements, the observable signature it should leave if the account is true. The interface is offered under a strict condition, stated before anything else, because the condition governs every sentence that follows.

\subsection{The condition: an interface is not an identity}
\label{p22:sub:neuro-condition}

The condition is the same one this series has imposed wherever it has reached for the language of a formal or empirical science, in the predictive reading of trust and in the dynamical reading of the relational field. The neuroscience is taken at the level of implementation, as the description of what should be observable if the higher-level account holds, and not as a reduction of the higher-level account to the mechanism. To say that reception leaves a signature in how prediction errors are integrated is not to say that reception is that integration, in the way that to say a vow leaves a signature in a signed document is not to say the vow is the ink. The relation asserted throughout this section is that the four movements, if actual, should leave observable signatures, and that relation is one of expected manifestation under a hypothesis, not of identity. No equals sign appears in this section, and its absence is not a stylistic tic but the exact mark of the commitment: an identity claim would say the aesthetic act simply is a neural event, which would dissolve the relational ontology the whole paper has built, whereas a signature claim says only that the act, being actual, should be found to have left a trace, which leaves the ontology intact and renders it answerable to evidence at the same time. The value of the interface is precisely that it makes the account falsifiable without making it reductive: if the signatures are sought and not found, the account is in trouble; if they are found, the account is supported but not thereby shown to be nothing but the mechanism.

\subsection{Perception: sensitivity to deviation}
\label{p22:sub:neuro-perception}

The perception section read the seeing of the field, at its implementation level, as a refusal to let an over-weighted prior flatten the field, a keeping of the weight on incoming evidence so that the unplanned deviation is not silently explained away. If that reading is right, then a perceiver who has cultivated this keeping should differ, observably, from one who has not, in responsiveness to deviation. The signature to look for is an enhanced neural response to the unexpected, the kind of response indexed by the established markers of surprise and salience, and a corresponding change after training: the prediction of the account is that the cultivation described in the perception section should, if it does what it claims, raise the system's registration of contingency rather than leaving it flat under a confident prior. The salience-marking systems that flag the unexpected, and the event-related components that rise to a deviation, are the candidate places such a signature would appear. The account does not assert that perceiving the field is the firing of these systems; it predicts that, if the cultivation of perception is real, its effect should be visible there, as a measurable shift in how much an unexpected event in the field is registered rather than absorbed in advance.

A subtlety must be added, because the naive form of this prediction is too crude and would be embarrassed by an obvious objection. If the signature of cultivated perception were simply a larger response to every deviation, then the cultivated perceiver would be one who is startled by everything, overwhelmed by surprise, unable to let the predictable be predictable, which is plainly not what a fine perception of the field is; a perception that responded maximally to every deviation would be not fine but pathological, drowning in noise. The prediction must therefore be refined: the signature of cultivated perception is not a uniformly larger response to deviation but a better-calibrated one, an enhanced response to the deviations that matter, the contingencies that bear on the field's state, and not to mere noise. The cultivated perceiver registers more strongly the unplanned thing that signifies, the remark that bears a weight, the shift that means something, while not being thrown by every random fluctuation, and this is a matter of what the system treats as signal rather than of how much it responds in general. The refined prediction is thus about the selectivity of the response, its tuning to the significant deviation, and not about its gross magnitude, and this selectivity is harder to measure but truer to what cultivated perception is, the discrimination of the deviation that matters from the noise that does not, which is exactly the resolution the coupling section described, now sought as a calibration of response rather than a mere amplitude of it. The naive prediction, larger responses, would be both easy to test and wrong; the refined prediction, better-calibrated responses, is harder to test and right, and the interface must commit to the harder and right one.

\subsection{Reception: integration rather than gating}
\label{p22:sub:neuro-reception}

The reception section drew a distinction between two fates of what perception delivers: it can be integrated, allowed to update the model and alter the valuation it enters, or it can be defensively gated, held off so that it changes nothing. That distinction, if real, has an implementation-level signature in whether a prediction error is propagated into an update or is suppressed before it can. The account predicts that reception, when it succeeds, should be visible as the integration of the unexpected into a revised model, the kind of revision that a measurable learning rate would register, and that defended reception should be visible as the suppression of that revision, an error met and gated rather than learned from. The regulation of limbic reactivity by prefrontal systems is the candidate locus for the difference between meeting an unexpected, possibly disturbing event with integration and meeting it with defence, and the developmental claim that injury raises defence and forecloses containment makes a further prediction here, that a relation or a subject under the load of recent injury should show, at this interface, more gating and less integration, a measurable correlate of the wounded gate the reception section described. Again the relation is one of signature and not identity: reception is not the prefrontal regulation of a limbic response, but if reception is actual it should leave its trace in whether the error it meets is integrated or suppressed.

This signature requires a refinement that distinguishes it from a confound the account must rule out, the confound of mere indifference. A system that integrated every error without suppression might look, at this interface, like perfect reception, but it might equally be a system that defended against nothing because nothing reached it, a flatness mistaken for openness, the indifference that lets everything in because it cares about none of it. Reception, as the paper means it, is not indifference; it is the letting-in of what matters, the integration of the significant error that defence would otherwise gate, and its signature must therefore distinguish the open reception of the significant from the indifferent passage of everything. The refined prediction is that reception shows as the integration of errors that carry weight, the unexpected things that bear on the field, while indifference shows as a uniform low reactivity that integrates because it is not engaged, and the difference between them is whether the integration is selective, tuned to what matters, or flat, equal across the significant and the trivial alike. This parallels the refinement the perception signature required, the distinction of the calibrated response from the merely large one, and it has the same source: reception, like perception, is a matter of selectivity, the letting-in of the significant and not the indiscriminate admission of everything, and its signature is the selective integration of what matters rather than the flat integration of all. The interface must therefore predict not integration as such but selective integration, the letting-in of the weighted error and not the indifferent passage of every error, and this selective signature is what distinguishes the reception the paper means from the indifference that would counterfeit it at the level of the mechanism.

The developmental dimension deepens this and connects it to the account's most humane claim. If reception is a capacity that injury damages, then the signature of damaged reception should be not a uniform suppression but a specific pattern, the gating of exactly the errors that bear on the wound, the defence raised against the significant while the trivial passes freely, the selective closure that injury produces. A subject whose reception has been wounded does not become indifferent to everything; she becomes defended against what touches the injury, integrating freely what is safe and gating what is dangerous, and the signature of her wounded reception is this selective gating, the closure tuned to the threatening rather than spread across all. This makes the interface's developmental prediction specific and testable: the wounded gate is not a global suppression but a targeted one, raised against the significant and dangerous, and its healing should show as the gradual return of integration to exactly the errors it had gated, the wound's specific closure specifically reopening as the injury heals. The mercy the reception section urged toward the wounded gate is thus given an implementation-level shape, the recognition that the wounded subject is not indifferent but specifically defended, and that her healing is the specific reopening of a specific closure, a process the interface predicts should be visible as the targeted return of integration to the errors the wound had taught her to gate.

\subsection{Sharing: inter-brain coupling}
\label{p22:sub:neuro-sharing}

The sharing interface is the one of greatest consequence, because it is the place where the central concept of the paper, relational aesthetics, is most directly exposed to evidence, and so it is worth stating with care what is and is not being claimed. The concept holds that beauty is a coupling event between subjects, produced in joint attention and not transmitted between two privacies. If that is so, then a shared aesthetic moment should differ from two simultaneous private ones in a way that is not located in either subject alone but in the relation between their two nervous systems, and the candidate signature is inter-brain synchrony, the coupling between two brains that the methods of dual recording, hyperscanning, are built to detect~\parencite{czeszumski_hyperscanning}. The prediction is specific: a genuinely shared aesthetic moment, the recursive common focus the sharing section described, should leave a signature in inter-brain coupling that the side-by-side coincidence of two private judgements should not, because the former has the between that the latter lacks, and the between is exactly what a measure spanning two brains, rather than residing in one, is positioned to register.

This makes inter-brain synchrony the strongest candidate the paper has for a hard empirical correlate of its central concept, and the stakes should be stated honestly in both directions. If shared aesthetic moments are found to be marked by inter-brain coupling that matched but unshared moments are not, that is direct empirical support for the claim that beauty in the relevant sense is a coupling event with a site between the subjects rather than within either, support of a kind that no analysis of single brains could provide. If, on the other hand, no such differential coupling is found, the concept is under real pressure, because a beauty claimed to be produced between two subjects ought to be visible between them, and its invisibility there would be evidence against the betweenness the concept asserts. The interface thus does for relational aesthetics what an interface should do: it stakes the concept on a finding rather than insulating it. And still the condition holds, that even a confirming synchrony would be the signature of the coupling and not the coupling itself, since the beauty of the shared moment is a relational event in the world and not a correlation in a recording, however much the correlation, if found, would be its trace.

A methodological objection must be faced, because it is the one the relevant research literature has learned to press, and an interface that ignored it would be naive. Inter-brain synchrony between two people attending to the same scene might be no coupling between them at all but a mere artefact of common input: two brains exposed to the same light, the same room, the same passing event, will show correlated activity simply because they are driven by the same external causes, and such correlation would mark no betweenness, only a shared outside. If the synchrony the account predicts were of this kind, it would not be evidence for relational aesthetics, since two strangers who happened to face the same scene would show it equally, and the concept requires a coupling between the two and not a common cause acting on both. The objection is serious, and the account does not wave it away; it specifies the prediction so as to exclude the artefact. What relational aesthetics predicts is not synchrony during shared attention as such but a differential synchrony, an excess of coupling in the genuinely shared aesthetic moment over the coupling present when the same two attend to the same scene without the recursive common focus that makes the moment shared in the relevant sense. The control is the matched but unshared moment, two people attending to the same light without the joint attention that constitutes the between, and the prediction is that the shared moment exceeds it. Common input is present in both the shared and the matched-unshared condition and so cannot explain a difference between them; only the coupling the concept names is present in the one and absent in the other. The prediction is therefore not the naive one that vulnerable to the artefact but the specific one that the artefact cannot produce, and stating it so is what makes inter-brain synchrony a genuine test of betweenness rather than a measure of shared illumination.

\subsection{Reproduction: experience-dependent plasticity}
\label{p22:sub:neuro-reproduction}

The reproduction section described accumulation as a lowered threshold of joint perception, a raised starting point carried across cycles, the settling of a changed disposition rather than a stored stock. The implementation-level signature of such accumulation is experience-dependent plasticity, the lasting change in a system wrought by repeated experience, expressed here as a downward shift in the threshold at which the partners jointly register what before went unseen. The prediction is that good reproduction, the spiral that gains, should be visible as a longitudinal change in perceptual sensitivity, a lowering across time of the threshold for joint registration, covarying with the frequency of shared aesthetic moments, whereas degenerate reproduction, the wheel, should show no such drift, the same gesture repeated leaving sensitivity flat. This is the slowest of the four signatures, the one that appears only over time, which is fitting, since reproduction is the movement that lives in time, and its trace is by nature a trace across cycles rather than within a moment. The non-hoarding character of the accumulation has its analogue here too: what plasticity records is not a stored quantity but a changed responsiveness, a disposition rather than a deposit, which is the implementation-level echo of the distinction between curvature and stock that the reception and reproduction sections drew.

\subsection{The purpose of the interface}
\label{p22:sub:neuro-purpose}

Gathered, the four signatures form a research programme more than a result, a set of places to look rather than a set of findings, and they are offered in that spirit. Their function in the paper is twofold. First, they discharge the debt the perception section incurred when it retained the predictive account at the level of implementation rather than ontology: the debt is paid by showing, across all four movements, exactly what implementation-level work the empirical sciences do here, which is to supply the observable signatures of an account whose substance lies elsewhere. Second, they keep the whole edifice answerable. A philosophy of intimacy that made claims about perception, reception, sharing, and reproduction while standing wholly beyond the reach of evidence would be asking to be believed rather than tested, and the series has consistently preferred to be tested. The interface is the form that preference takes here: it names, for each movement, the finding that would support the account and the absence that would embarrass it, and it does so without surrendering the relational ontology to the mechanism, holding throughout to the one discipline this section exists to keep, that a signature is the trace of a thing and never the thing itself.

\subsection{The interface and the relational ontology}
\label{p22:sub:neuro-ontology}

A final word situates the whole interface within the relational ontology it serves, lest the gathering of neural signatures be misread as a quiet concession that the relational account was, after all, reducible to the brain. The opposite is the case, and the structure of the interface shows why. Of the four signatures, the one that bears most directly on the paper's central concept, the inter-brain coupling that would mark the sharing in which beauty is produced, is precisely the one that cannot be located in any single brain, that requires a measure spanning two nervous systems, that is by its nature a signature of the between and not of either pole. The most important signature the interface seeks is therefore one that a single-brain neuroscience could never find, a coupling that exists only between two brains and not within either, and this is the strongest possible indication that the interface, far from reducing the relational to the neural, requires a neuroscience that has itself become relational, that measures couplings between brains and not only states within them. The relational ontology is not conceded by the interface but imposed upon the neuroscience by it, the concept of relational aesthetics demanding a measure that spans the between because the beauty it names lives there.

This is why the interface, for all its empirical commitment, leaves the relational ontology not weakened but confirmed in its place. The signatures are the traces the relational events leave in the implementing machinery, and the most central of them is a trace that only a between can leave, a coupling no single brain contains. To seek it is to commit the neuroscience to the very betweenness the paper's ontology asserts, to require of the measurement exactly the relational structure the concept describes, so that the interface, rather than dissolving the relational into the neural, extends the relational into the neural, carrying the betweenness of the concept into the very design of the measurement that would test it. The implementation level, properly understood, does not undercut the ontology above it but answers to it, shaped by the relational structure of what it implements, and the interface is therefore the relational ontology's reach into the empirical, not its surrender to it, the point at which a philosophy of the between requires and receives a science of the between, each holding to the relational structure that the other, in its own register, describes.
